%%%%%%%%%%%%%%%%%%%%%%%%%%%%%%%%%%%%%%%%%%%%%%%%%%%%%%%%%%%%%%%%%%%%%%%%%
%                                                                       %
%                          Szablon                                      %
%                                                                       %
%%%%%%%%%%%%%%%%%%%%%%%%%%%%%%%%%%%%%%%%%%%%%%%%%%%%%%%%%%%%%%%%%%%%%%%%%
%\documentclass[12pt]{article}
%%% aby uzyc amsart w pelnej krasie:
% 1. odkomentowac \address i \email
% 2. odkomentowac \keywords{} \subjclass
\documentclass[12pt]{amsart}
\usepackage{amssymb} % ten pakiet jest nieodzowny m.in. po to aby zrenderowac instrukcje varsubsetneq
\usepackage{amsmath}
\usepackage{latexsym}
\usepackage{amsfonts}
\usepackage{enumerate}
%\usepackage[mathscr]{eucal}
\usepackage[mathscr]{euscript}
\usepackage{eqlist}
\usepackage{amsthm}
\usepackage{mathtools}
\usepackage{enumitem} % po to aby mozna bylo tworzyc otoczenie
% enumerate z innymi sposobami numeracji itemow
\usepackage{hyperref} % aby moc tworzyc linki w tekscie
%%%%%%%%%%%%%%% Polish letter packages %%%%%%%%%%%%%%%%%%%%%%%%%%%%%%%%
%\usepackage[polish]{babel}
%\usepackage[utf8]{inputenc}
%\usepackage{t1enc}
%%% pakiet do generowania belkotu (wypelniacza):
\usepackage{lipsum}
%%% pakiet do dodawania fikusnych notek `a la ``todo''
\usepackage{todonotes}
%Ten pakiet niby jest m.in. po to aby mozna bylo ,,lamac'' linie w komorkach i ,,glowie'' tabeli
% (otoczenie table)
\usepackage{makecell}
%ten pakiet sluzy do eleganckiego formatowania listingow (programow? lecz zapewne nie tylko)
\usepackage{listings}
%ten pakiet sluzy do sprawdzania m.in. czy odpowiednie pozycje bibliograficzne
%sa faktycznie uzywane w tresci pliku. Przypuszczalnie nalezy go
%uruchamiac tylko podczas weryfikacji pliku - w oficjalnej wersji powininen byc wylaczony.
\usepackage{refcheck}
%pakiet dolaczony na razie po to aby moc utworzyc ,,monstrualne'' kwantyfikatory bigforall i bigexists
\usepackage{graphicx}
%pakiet rotating służy do obracania obiektów, takich jak tekst, tabele, obrazy, wykresy itp.
\usepackage{rotating}
%pakiet do zrecznego tworzenia wielowierszowych tabel.
\usepackage{multirow}
% nie dość że zawiera wiele symboli matematycznych to jeszcze daje
% lepszy (bardziej profesjonalny lub wyrazisty) design.
\usepackage{stix}
% pierwszy pakiet daje możliwość użycia fontów caligraficznych,
% natiomiast nie mam pojęcia w jakim celu jest ten drugi pakiet,
% ale on jest z całą pewnością powiązany z pierwszym.
\usepackage{calligra}
\usepackage[T1]{fontenc}
%%% to jest celem podkreślenia tekstu przez ,,wężyk''.
\usepackage{ulem}
\normalem %aby przywrócić standardowe funkcjonowanie \emph.
% po to by używać bibera:
\usepackage[backend=biber]{biblatex}

%%% to pozwala kontrolować overfulle, są one zaznaczane na marginesie czarną beleczką (o zadanym rozmiarze):
\overfullrule=5pt 

%---------------------------------------------------------------------
% Theorems 
%---------------------------------------------------------------------
 
% Theorem style 'plain' are for: Theorem, Lemma, Corollary,
% Proposition, Conjecture, Criterion, Algorithm
\theoremstyle{plain}
\newtheorem{theorem}{Theorem}[section]
\newtheorem{lemma}[theorem]{Lemma}
\newtheorem{corollary}[theorem]{Corollary}
\newtheorem{conclusion}[theorem]{Corollary}
\newtheorem{claim}[theorem]{Claim}
\newtheorem{fact}[theorem]{Fact}
\newtheorem{proposition}[theorem]{Proposition}
\newtheorem{axiom}{Axiom}
% Theorem style 'definition' are for: Definition, Condition, Problem,
% Example
\theoremstyle{definition}
\newtheorem{definition}[theorem]{Definition}
\newtheorem{example}[theorem]{Example}
\newtheorem{exercise}{Exercise}
\newtheorem*{solution}{Solution}
\newtheorem{remark}[theorem]{Remark}
\newtheorem{Problem}[theorem]{Problem}
% Theorem style 'remark' are for: Remark, Note, Notation, Claim,
% Summary, Acknowledgement, Case, Conclusion
\theoremstyle{remark}
\newtheorem*{notation}{Notation}
\newtheorem*{acknowledgement}{Acknowledgement}

%%%%%% makra:

%A
\newcommand{\afc}{\mathrm{AFC}}
\newcommand{\afcbar}{\overline{\afc}}
\newcommand{\arr}{\rightarrow}
\newcommand{\Arr}{\Rightarrow}

%B
\newcommand{\baire}{\omega^{\omega}}
\newcommand{\Bor}{\mbox{$\mathcal{Bor}$}}

%C
\newcommand{\ca}{2^{\omega}}
\newcommand{\cantor}{\ca}
\newcommand{\Card}[1]{\Vert #1 \Vert}

%D
\newcommand{\dom}{{\rm dom}}
\newcommand{\dummy}{{\tt Blah blah blah}}

%E
\newcommand{\Even}{\hbox{\rm \tiny Even}}

%F
\newcommand{\finsub}{[\omega]^{<\omega}}
\newcommand{\forces}{\mathrel{\|}\joinrel\mathrel{-}}

%G
\newcommand{\Graph}{\hbox{\it Graph}}
\newcommand{\Graphit}{{\mathit Graph}}
\newcommand{\Graphh}{\operatorname{Graph}}

%H
\newcommand{\homeomorphic}{\approx}

%I
\newcommand{\incr}{\omega^{\uparrow \omega }}
\newcommand{\infsub}{[\omega]^{\omega}}

%L
\newcommand{\la}{\langle}

%M
\newcommand{\meager}{{\cal MGR}}
\newcommand{\minideal}{${\cal F}_{\hbox{\rm \scriptsize min}}(\neg
  D)\;$}

%N
\newcommand{\neglig}{{\cal N}}
\newcommand{\nnatural}{\mathbb{N}}

%O
\newcommand{\Odd}{\hbox{\rm \tiny Odd}}

%P
\newcommand{\Part}{{\it Part}}
\newcommand{\Perf}{{\it Perf}}
%%%\newcommand{\proof}{\flushleft{ \sc Proof. } \\ }
\newcommand{\Proof}[1]{\bigbreak\noindent{\bf Proof #1}\enspace}

%Q
%%%\newcommand{\qed}{{\hfill\vrule height6pt width6pt depth1pt\medskip}}
%\newcommand{\qed}{\sharp}
% to poniżej w komentarz o ile używamy pakietu stix
%\newcommand{\QED}{\hspace{0.1in} \Box \vspace{0.1in}}

%R
\newcommand{\ra}{\rangle}
\newcommand{\ran}{{\rm ran}}
\newcommand{\rational}{\mathbb{Q}}
\newcommand{\real}{\mathbb{R}}

%S
\newcommand{\seq}{\subseteq}
%%%\newcommand{\square}{\hbox{\ \ \ \ \ \vrule\vbox{\hrule\phantom{o}\hrule}\vrule}}
% a small restriction:
\newcommand{\srestriction}{{\hbox{${\scriptstyle\,|\grave{}\,}$}}}
\newcommand{\supp}{\mathit{supp}}

%U
\newcommand{\up}{\uparrow}
\newcommand{\ucrz}{UCR_0}

%%%%%%%%%%%%%%%%%%%%%% Calligraphic font commands %%%%%%%%%%%%%%%%%%%%%%%%%%
\newcommand{\calA}{\mathcal{A}}
\newcommand{\calB}{\mathcal{B}}
\newcommand{\calC}{\mathcal{C}}
\newcommand{\calD}{\mathcal{D}}
\newcommand{\calE}{\mathcal{E}}
\newcommand{\calF}{\mathcal{F}}
\newcommand{\calG}{\mathcal{G}}
\newcommand{\calH}{\mathcal{H}}
\newcommand{\calI}{\mathcal{I}}
\newcommand{\calJ}{\mathcal{J}}
\newcommand{\calK}{\mathcal{K}}
\newcommand{\calL}{\mathcal{L}}
\newcommand{\calM}{\mathcal{M}}
\newcommand{\calN}{\mathcal{N}}
\newcommand{\calO}{\mathcal{O}}
\newcommand{\calP}{\mathcal{P}}
\newcommand{\calQ}{\mathcal{Q}}
\newcommand{\calR}{\mathcal{R}}
\newcommand{\calS}{\mathcal{S}}
\newcommand{\calT}{\mathcal{T}}
\newcommand{\calU}{\mathcal{U}}
\newcommand{\calV}{\mathcal{V}}
\newcommand{\calW}{\mathcal{W}}
\newcommand{\calX}{\mathcal{X}}
\newcommand{\calY}{\mathcal{Y}}
\newcommand{\calZ}{\mathcal{Z}}
%%%%%%%%%%%%%%%%%%%%%% inne %%%%%%%%%%%%%%%%%%%%%%%%%%%%%%%%%%%%%%%
\newcommand{\SqrFr}{\mathbb{SF}}
\newcommand{\Primes}{\mathit{Primes}}
\newcommand{\mathint}{\mathit{int}}
\newcommand{\mathcl}{\mathit{cl}}
\newcommand{\exampleGFnotK}{\mathit{Ex}}
%%%%%%%%%%%%%%%%%%%%%% Some Greek fonts %%%%%%%%%%%%%%%%%%%%%%%%%%%%%%%%%%%%%%%
\newcommand{\oo}{\omega}
\newcommand{\bb}{\beta}
\newcommand{\dd}{\delta}
\newcommand{\ee}{\varepsilon}
\newcommand{\kk}{\kappa}
%%%\newcommand{\th}{\theta}
\newcommand{\stirlingii}{\genfrac{\{}{\}}{0pt}{}}
%%% makra dla utworzenia monstrualnych kwantyfikatorow:
\newcommand{\bigforall}{\mbox{\Large $\mathsurround0pt\forall$}}
\newcommand{\bigexists}{\mbox{\Large $\mathsurround0pt\exists$}}
%%% tutaj makra specyficzne tylko w tej pracy:
\newcommand{\NICEDIRECTIONS}{\mathrm{NiceDir}}

%\renewcommand\abstractname{Summary}
%%%To print the date and time on each page
%%% comment out if not needed (next 14 lines)
\makeatletter
{\newcount\@hour}
{\newcount\@minute}
\def\timenow{\@hour=\time \divide\@hour by 60
\number\@hour:
  \multiply\@hour by 60 \@minute=\time
  \global\advance\@minute by -\@hour
  \ifnum\@minute<10 0\number\@minute\else
  \number\@minute\fi}
\def\ctimenow{\hfil{\tt \jobname.tex, \today~Time: \timenow }\hfil}
      \let\@oddfoot\ctimenow\let\@evenfoot\ctimenow
\makeatother

\pagestyle{myheadings}
\markboth{{\bf Andrzej Nowik}}{\bf \today}

%%%% opcjonalne dolaczanie komentarzy na marginesie (uruchomic zdejmujac komentarze z kodu ponizej):
%\typein[\wersjaRELEASE]{Czy wygenerowac wersje RELEASE ? (t/n)}
%\if\wersjaRELEASE t
%\typeout{Zostanie wygenerowana wersja RELEASE}
%\else
%\typeout{Zostanie wygenerowana wersja ROBOCZA}
%\fi
%\if\wersjaRELEASE t
%\def\uwzglednicKomentarze{n}
%\else
%\def\uwzglednicKomentarze{t}
%\fi
%\if\uwzglednicKomentarze t
%\typeout{Do kompilacji ZOSTANA dolaczone komentarze}
%\else
%\typeout{Plik przekompiluje sie BEZ dolaczania komentarzy}
%\fi
%%%% jesli powyzej jest w komentarzu to na stale dodajemy opisy na marginesie:
\def\uwzglednicKomentarze{t}
\newcommand{\comentario}[1]{\if\uwzglednicKomentarze t \marginpar{\tt #1} \fi}
% To see corrections comment next line and uncomment the second one
\newcommand{\correction}[2]{#1}
% \newcommand{\correction}[2]{#2}

\pagestyle{myheadings}
% To see corrections comment next line and uncomment the second one


\author{Andrzej Nowik}
\address[]{
Andrzej Nowik\\
University of Gda\'nsk \\
Institute of Mathematics \\
Wita Stwosza 57,
80--952 Gda\'nsk, Poland \\
}
\email[A.~Nowik]{andrzej.nowik@ug.edu.pl}

\addbibresource{rotacje_plik_pomocniczo_tymczasowy.bib}
\begin{document}
\title[
% a brief title
On rotations of sets...
]{
%a full title
On rotations of sets of measure zero and sets of first category
}
\keywords{topologia gestosci, topologia dyskretna, idealy}
\subjclass[2010]{03E05, 54D05, 54H05, 54A05, 11A41, 11B05, 11B25}

\begin{abstract}
dummy
\end{abstract}

\maketitle

\section{Notations}
Let $\mu$ and $\mu_2$ denote the standard one-dimensional Lebesgue measure 
and the two-dimensional Lebesgue measure, respectively.
Let $\theta \in [0, \pi)$ be any angle.
Let $\phi_\theta \colon \real^2\to\real^2$ denote
the rotation of the plane about the origin by an angle $\theta$.
Suppose that $B\subseteq\real$ is
a Lebesgue measurable set.
Let $\Phi(B)$ denote the set of all Lebesgue density points of the set $B$.

Let $\pi_x$ and $\pi_y$ denote the projections of the plane
onto the $x$-axis and $y$-axis, respectively.

If $Y$ is any subset of a topological space $X$ then by a portion of
$Y$ we mean a nonempty and relative open subset of $Y$
(see \cite{Laczkovich1998}, page 1783).

Under ``fat'' perfect set we mean a perfect set such that each of
its portion is of positive Lebesgue measure (see \cite{Buczolich2004}, page 497).
Let us note here that in the cited paper by Buczolich,
the notion of ``fat'' applies to Cantor sets
(i.e., perfect sets homeomorphic to the Cantor space),
but there is no obstacle to defining this property for all perfect sets.
\section{Preliminaries}

Let us recall the classical notion of a density point in the plane
(see for example \cite{WagnerBojakowska19921993}):

\begin{definition}
Suppose that $E \subseteq \real^2$ is a measurable set. 
Define
\[
  (x_0, y_0) \in \Phi_2(E) \iff
  \lim_{\varepsilon\to 0^+}
  \frac{\mu_2(E \cap [x_0 - \varepsilon; x_0 + \varepsilon] \times
   [y_0 - \varepsilon; y_0 + \varepsilon])}{4 \varepsilon^2} = 1.
\]
\end{definition}

It is well known that:
\begin{fact}
  This definition is equivalent to the ``circular'' one, namely:
\[
  (x_0, y_0) \in \Phi_2(E) \iff
  \lim_{\varepsilon\to 0^+}
  \frac{\mu_2(E \cap B((x_0, y_0), \varepsilon))}{\pi \varepsilon^2} = 1.
\]  
\end{fact}

\begin{claim}\label{claim:density-point-rotation}
The latter version of the definition of the $\Phi_2$ points is invariant under rotations,
so we have
\[
  (x_0, y_0) \in \Phi_2(E) \iff \phi_\theta(x_0, y_0) \in \Phi_2(\phi_\theta[E]).
\]
\end{claim}
\begin{lemma}\label{lemma:rotation-measurable-stripe}
  Suppose that $B \subseteq \real$ is a Lebesgue measurable set.
  Then $\Phi(B) \times \real \subseteq \Phi_2(B \times \real)$.
  \todo{Chyba tu zachodzi nawet równość?}
\end{lemma}
\begin{proof}
  Suppose $(x_0, y_0) \in \Phi(B)$ and let $0 < \eta < 1$.
  There exists $r_0 > 0$ such that
  \[
\forall_{0 < r < r_0} 
  \frac{\mu(B \cap [x_0 - r; x_0 + r])}{2r} > 1 - \eta.
\]
Then
\[
  \frac{\mu_2((B \times \real) \cap [x_0 - r; x_0 + r] \times
   [y_0 - r; y_0 + r])}{4 r^2} = 
\]
\[
  = \frac{\mu(B \cap [x_0 - r; x_0 + r]) \cdot 2r}{4r^2} > 1 - \eta.
\]
\end{proof}

\begin{lemma}
  For $\theta \in [0, \pi)$, $B\subseteq\real$ Lebesgue measurable, then
  $\phi_{\theta} [\Phi(B) \times \real] \subseteq \Phi_2(\phi_{\theta} [B\times\real])$.
\end{lemma}

\begin{proof}
  This is a consequence of the fact that density points in the plane are
  invariant under rotations (see Claim \ref{claim:density-point-rotation}).
  By Lemma \ref{lemma:rotation-measurable-stripe} we have:
  $\phi_{\theta} [\Phi(B) \times \real] \subseteq \phi_{\theta}[\phi_2(B\times\real)] =
  \Phi_2(\phi_\theta [B\times\real])$
\end{proof}

\begin{lemma}
  For $\theta_1,\ldots,\theta_L \in [0; \pi)$
  and for $B\subseteq\real$ Lebesgue measurable,
  we have
  \[
    \bigcap_{j=1}^L \phi_{\theta_j} [\Phi(B) \times \real] \subseteq
    \Phi_2(\bigcap_{j=1}^L \phi_{\theta_j} [B\times\real]).
  \]
\end{lemma}
\begin{proof}
  We have
  $\bigcap_{j=1}^L \phi_{\theta_j} [\Phi(B) \times \real] \subseteq$
  $\bigcap_{j=1}^L \Phi_2(\phi_{\theta_j} [B\times\real]) \subseteq$
  $\Phi_2(\bigcap_{j=1}^L \phi_{\theta_j} [B\times\real])$.
\todo{(...)}
\end{proof}

Let us define
\begin{definition}
  Suppose that $P \subseteq \real^2$ is a perfect plane set.
  Suppose that $\theta \in [0, \pi)$. We say that
  $P$ is {\it nowhere dense nice in the direction $\theta$}
  if, for every nowhere dense set $E \subseteq \real$, the set
  $P \cap \phi_{\theta}[E \times \real]$ is
  nowhere dense in the subspace topology on $P$.
  For a perfect set $P \subseteq \real^2$
  let $\NICEDIRECTIONS(P)$ denote the set of all directions $\theta \in [0, \pi)$
  such that $P$ is nowhere dense nice in the direction $\theta$.
\end{definition}
Let us prove the following easy characterization:
\begin{lemma}
Suppose that $P \subseteq \real^2$ is a perfect plane set and
suppose that $\theta \in [0, \pi)$. The following conditions are
equivalent:

\begin{enumerate}
\item
  $\theta \in \NICEDIRECTIONS(P)$;
\item
  For any open set $W \subseteq \real^2$ such that
  $W \cap P \not= \emptyset$ we have
  $\pi_x[\phi_{\theta}^{-1}[W \cap P]$ has a nonempty interior.
\end{enumerate}
 

\end{lemma}

\section{Results}

\begin{theorem}
  \label{theorem:main}
  Suppose that $(\mu_j)_{j=0}^\infty$ and $\theta_0,\ldots, \theta_{k-1}$
  are sequences of angles from $[0,\pi)$. Assume also that
  $\forall_{j \in 0,1,\ldots} \forall_{l = 0,\ldots, k-1} \mu_j \not= \theta_l$.
  Suppose also that $M \subseteq \real$ is a meager set and
  $N \subseteq \real$ is a Lebesgue measure zero set.
  Then the set
  \[
    \bigcup_{j = 0}^\infty \phi_{\mu_j}[M \times \real] \cup
    \bigcup_{l = 0}^{k-1} \phi_{\theta_l}[N \times \real]
  \]
  is not the whole plane $\real^2$.
\end{theorem}

\begin{proof}
  Suppose, for contradiction, that the conclusion is false.
  Without loss of generality, we may assume that
  $M$ is an $F_{\sigma}$ meager set and $N$ is a $G_{\delta}$
  Lebesgue null set.

  First, we will prove the following lemma:

  \begin{lemma}\label{lemma:main}
  Suppose that $\theta_0,\ldots,\theta_{k-1}$
  is a finite sequence of angles from $[0,\pi)$, and suppose that 
  $N \subseteq \real$ is a Lebesgue measure zero set.
  Then there exist a perfect set $P$ such that
  \[
     P \cap \bigcup_{l = 0}^{k-1} \phi_{\theta_l}[N \times \real] = \emptyset.
  \]
  and, moreover,
    $[0, \pi) \setminus \lbrace \theta_0,\ldots,\theta_{k-1} \rbrace \subseteq \NICEDIRECTIONS(P)$.
  \end{lemma}
  
\begin{proof}
  Choose a closed set $E \subseteq \real\setminus N$
  such that $\overline{\Phi(E)} = E$ (thus $E$ is a perfect set)
  and moreover, all portions of $E$ have measure greater than zero.
  We may also assume that
  \[
    \mu_2(\bigcap_{l = 0}^{k-1} \phi_{\theta_l}[E \times \real]) > 0.
  \]
  The latter follows from the following identity:
  For any increasing sequence of sets: $E_n \subseteq E_{n+1}$
  we have
  \[
    \bigcap_{l = 0}^{k-1} \phi_{\theta_l}[\bigcup_{n=1}^\infty E_n \times \real]
    = \bigcup_{n=1}^\infty
      \bigcap_{l = 0}^{k-1} \phi_{\theta_l}[E_n \times \real].
  \]

  Let us denote:
  \[
   E^* = \bigcap_{l = 0}^{k-1} \phi_{\theta_l}[E \times \real]
  \]
  \[
  \tilde{E} = \bigcap_{l = 0}^{k-1} \phi_{\theta_l}[\Phi(E) \times \real].
\]
and \[
P^* = \overline{\Phi_2(E^*)}.
\]
We have
\[
\tilde{E} \subseteq \Phi_2(E^*) \subseteq P^* \subseteq E^*.
\]
Moreover, $P^*$ is a perfect ``fat'' set.  
It suffices to check only that this set is a ``fat'' set,
i.e. if $W$ is an open set such that $W \cap \overline{\Phi_2(E^*)} \not= \emptyset$
then $\mu_2(W \cap \overline{\Phi_2(E^*)}) > 0$. But since
$W \cap \overline{\Phi_2(E^*)} \not= \emptyset$, then
$W \cap \Phi_2(E^*) \not= \emptyset$, hence
$\mu_2(E^* \cap W) > 0$.

We have also
$\mu_2(E^* \setminus \Phi_2(E^*)) = 0$, thus
$\mu_2(\Phi_2(E^*) \cap W) > 0$.

Let us observe that
\[
  E^* \setminus \tilde{E} =\]
\[
  \bigcap_{l = 0}^{k-1} \phi_{\theta_l}[E \times \real] \setminus
  \bigcap_{l = 0}^{k-1} \phi_{\theta_l}[\Phi(E) \times \real] \subseteq
\]
\[
   \bigcup_{l = 0}^{k-1} \phi_{\theta_l}[(E \setminus \Phi(E)) \times \real]
\]
and this latter set has two-dimensional Lebesgue measure zero.

Observe that we have
\[
  E^* \cap \bigcup_{l = 0}^{k-1} \phi_{\theta_l}[N \times \real] = \emptyset.
\]
Indeed, suppose that $X$ is a point from $E^*$ such that
$X \in \bigcup_{l = 0}^{k-1} \phi_{\theta_l}[N \times \real]$,
so there exists $0 \leq l_0 \leq k - 1$ and $(x(L), y(L)) \in N \times \real$
such that $X = \phi_{\theta_{l_0}}((x(L), y(L)))$.
Since $X \in E^*$, $X \in \phi_{\theta_{l_0}}[E \times \real]$.
Hence $(x(L), y(L)) \in E \times \real$, which is a contradiction with
the definition of the set $E$.

By way of contradiction, suppose that
$\mu \in [0, \pi) \setminus \lbrace \theta_0,\ldots,\theta_{k-1} \rbrace$
and there exists a nowhere dense closed set $F \subseteq \real$
such that $\phi_{\mu}[F \times \real] \cap P^*$ is not a nowhere dense
set in the subspace topology on $P^*$.

Then there exists an open $W^* \subseteq \real^2$ 
such that $\emptyset \not= W^* \cap P^* \subseteq \phi_{\mu}[M \times \real]$.
Hence $W^* \cap \Phi_2(E^*) \not= \emptyset$, therefore
$\mu_2(W^* \cap E^*) > 0$. Since $\mu_2(E^* \setminus \tilde{E}) = 0$
we obtain that $W^* \cap \tilde{E} \not= \emptyset$.

Take any point $X_0 \in W^* \cap \tilde{E}$. So we have
\[
  X_0 \in W^* \cap \bigcap_{l = 0}^{k-1} \phi_{\theta_l}[\Phi(E) \times \real].
\]
Denote:
\begin{align*}
\phi^{-1}_{\mu}(X_0) &= (x(\mu), y(\mu)) \\
\phi^{-1}_{\theta_l}(X_0) &= (x^{(l)}_0, y^{(l)}_0) \qquad \text{for }l = 0,\ldots, k-1 
\end{align*}

Evidently
$x(\mu) \in M$ and $x^{(l)}_0 \in \Phi(E)$ for $l = 0, \ldots, k-1$.

Let us formulate the following auxiliary lemma:
\begin{lemma}
  If $B \subseteq A$, $A, B$ Lebesgue measurable sets, $\mu(B) > 0$, $\mu(A) < \infty$,
  $E$ also Lebesgue measurable set, then if
  for some $\varepsilon \in (0,1)$ we have
  \[
  \mu(E \cap A) \geq (1 - \varepsilon) \mu(A),
\]
then
\[
  \mu(E \cap B) \geq \big(1 - \varepsilon \frac{\mu(A)}{\mu(B)}\big) \mu(B),
\]
\end{lemma}
\begin{proof}
  We have
  $\mu(E \cap B) = \mu(E \cap A) - \mu(E \cap A \setminus B) \geq$
  $(1 - \varepsilon) \mu(A) - \mu(A \setminus B) =
  (1 - \varepsilon) \mu(A) - \mu(A) + \mu(B) =$
  $- \varepsilon \mu(A) + \mu(B) = $
  $\mu(B) \big(1 - \varepsilon \frac{\mu(A)}{\mu(B)}\big)$.
\end{proof}

For $l = 0,\ldots, k-1$
let us denote the angle $\mu - \theta_l$ by $\theta^*_l$.
By our assumption, $\theta^*_l \not= 0$, hence
$\theta^*_l \in (-\pi, \pi) \setminus \lbrace 0 \rbrace$.
We have
$\pi_x(\phi_{\theta_l}^{-1}(\phi_{\mu}(x(\mu), y(\mu))) =$
$ \pi_x(\phi_{\theta_l}(X_0)) =$
$\pi_x(x^{(l)}_0, y^{(l)}_0) = x^{(l)}_0$.
Likewise:
$\pi_x(\phi_{\theta_l}^{-1}(\phi_{\mu}(x(\mu), y(\mu))) =$
$\pi_x(\phi_{\theta^*_l}(x(\mu), y(\mu)) =$
$x(\mu) \cos \theta^*_l - y(\mu) \sin \theta^*_l$.

Let us choose $\varepsilon > 0$ such that
\[
  \frac{1}{4 \sum_{l=0}^{k-1} \frac{1}{|\sin \theta^*_l}|} > \varepsilon.
\]
and let $\eta_0 > 0$ be such that
\begin{enumerate}
\item  
For each $\hat{x}(\mu), \hat{y}(\mu)$ such that
$|\hat{x}(\mu) - x(\mu)| < \eta_0$ and
$|\hat{y}(\mu) - y(\mu)| < \eta_0$, we have
$\phi_{\mu}(\hat{x}(\mu), \hat{y}(\mu)) \in W^*$.
%%% this is true since $\phi_{\mu}(x(\mu), y(\mu)) = X_0 \in W^*$.
\item
  For each $0 < \eta < \eta_0$ we have
$\mu((x^{(l)}_0 - \eta, x^{(l)}_0 + \eta$
\end{enumerate}

\end{proof}

Continuation of the proof of Theorem \ref{theorem:main}:

Apply Lemma \ref{lemma:main} to the sequence $\theta_0,\ldots,\theta_{k-1}$
and to the set $N$. So we get a suitable perfect set $P$.
Observe that we have
\begin{equation}\label{equation:P-is-pathological}
  P \subseteq \bigcup_{j = 0}^\infty \phi_{\mu_j}[M \times \real].
\end{equation}
Indeed, suppose that $X$ is a point from $P$ such that
$X \not\in \bigcup_{j = 0}^\infty \phi_{\mu_j}[M \times \real]$.
Then $X \in \bigcup_{l = 0}^{k-1} \phi_{\theta_l}[N \times \real]$,
which contradicts the statement of Lemma \ref{lemma:main}.
Since $\mu_j \in \NICEDIRECTIONS(P)$ for all $j = 0,1,\ldots$, 
$\phi_{\mu_j}[M \times \real]$ is a meager set in the
subspace topology on $P$. By the inequality (\ref{equation:P-is-pathological})
we obtain a contradiction.

\end{proof}
%%%%%%%%%%%%%%%%%%%%%%%%%%%%%%%%%%%%%%%%%%%%%%%%%%%%%%%%%%%%%%%%%%%%%%%%%%%%% 
%                        The Bibliography                                   %
%%%%%%%%%%%%%%%%%%%%%%%%%%%%%%%%%%%%%%%%%%%%%%%%%%%%%%%%%%%%%%%%%%%%%%%%%%%%%

\begin{filecontents*}[overwrite]{rotacje_plik_pomocniczo_tymczasowy.bib}
@article {WagnerBojakowska19921993,
    AUTHOR = {Wagner-Bojakowska, El\.zbieta},
     TITLE = {Some remarks on density topologies on the plane},
   JOURNAL = {Real Anal. Exchange},
  FJOURNAL = {Real Analysis Exchange},
    VOLUME = {18},
      date = {1992/1993},
    NUMBER = {2},
     PAGES = {339--342}
  %    ISSN = {0147-1937,1930-1219},
  % MRCLASS = {26A03 (54A10 54A25 54C08)},
  %MRNUMBER = {1228399},
  %MRREVIEWER = {Aliasghar\ Alikhani-Koopaei},
}
@article{Laczkovich1998,
  author  = {M. Laczkovich},
  title   = {Analytic subgroups of the reals},
  journal = {Proc. Amer. Math. Soc.},
  volume  = {126},
  year    = {1998},
  pages   = {1783--1790}
}
@article{Buczolich2004,
  author  = {Zolt{\'a}n Buczolich},
  title   = {Category of density points of fat Cantor sets},
  journal = {Real Analysis Exchange},
  volume  = {29},
  number  = {1},
  pages   = {497--502},
  date    = {2003/2004}
}
\end{filecontents*}

\printbibliography

\begin{center}
\ctimenow
\end{center}

\end{document}
